\documentclass[12pt]{article}
\usepackage[utf8]{inputenc}
\usepackage[T1]{fontenc}
\usepackage{hyperref}
\title{Sinicahua Mixtec Language Resources for K-12 Curriculum Development}
\author{Compiled by Claude Sonnet 5.5 (claude-sonnet-5-5) for Andr\'as Kornai}
\date{}
\begin{document}
\maketitle

\section{Introduction}
Sinicahua Mixtec (ISO 639-3 xti; Glottocode sini1243) is an Otomanguean language of Oaxaca, Mexico, in the Mixtec subgroup. Glottolog places it under Chalcatongic and Sinicahua-Tijaltepec. Alternative names include Mixteco de San Antonio Sinicahua, Mixteco de Yoloxochitl and Mixteco: Siniyucu. The searches found only a few resources, and no grammars, dictionaries or school materials.

\section{Primary Educational Resources}
None found. The language is not known to be taught in schools.

\section{Linguistic and Reference Materials}
\begin{itemize}
\item Glottolog entry for Sinicahua Mixtec, with classification and a reference list worth following up. \url{https://glottolog.org/resource/languoid/id/sini1243}
\end{itemize}
No grammar or dictionary specific to this variety turned up. Resources for other Mixtec varieties should not be assumed to apply.

\section{Religious and Community Resources}
Global Recordings Network offers Good News audio-visual Bible lessons (40 sections with pictures) for evangelism and church planting. They are religious in nature and would need community review before any school use. \url{https://www.globalrecordings.net/en/language/xti}

\section{Digital and Audio Resources}
The Global Recordings page provides MP3 audio, low-bitrate audio and MP4 slideshow downloads, which could serve as native-language listening samples subject to community consent. \url{https://www.globalrecordings.net/en/language/xti}

\section{Curriculum Development Context}
Sinicahua Mixtec is reported to be used as a first language by the community, which is a strong basis for native-language instruction. Almost no written material was found. A curriculum would need to start with community-led work: orthography, oral story collection, and Spanish-Mixtec bilingual materials. Related Mixtec resources could inform the work but need adaptation for this variety.
\end{document}